\setcounter{zagc}{17}\zag{Zasada domniemania niewinności oraz zasada in dubio pro reo}[{\kolor A}]\str{201-209}

Zasada domniemania niewinności w polskim systemie prawnym oparta jest na art. 42 ust. 3 \acr{krp}, który, adresowany do państwa, brzmieniem „\textit{Każdego uważa się za niewinnego, dopóki jego wina nie zostanie stwierdzona prawomocnym wyrokiem sądu.}” nakazuje wszystkim organom państwowym traktować każdego człowieka, a zwłaszcza podejrzanego i oskarżonego, jak niewinnego i zakładać jego niewinność aż do stwierdzenia uprawomocnionym wyrokiem sądu, że niewinny nie jest.
\Acr{kpk} w art. 5 § 1 ogranicza się do objęcia domniemaniem tylko oskarżonego.\medskip

Domniemanie to nakłada ciężar dowodu winy oskarżonego na oskarżyciela i niedopuszczalne jest przenoszenie tego ciężaru na obronę, co wielokrotnie potwierdzał w swoich wyrokach \acr{etpcz}. Jest ono ściśle powiązane z prawomocnością orzeczenia i chroni oskarżonego aż do uprawomocnienia a w razie uchylenia przełamującego je prawomocnego wyroku odżywa.\medskip

Nie stoi domniemanie niewinności na przeszkodzie stosowania środków przymusu procesowego, albowiem, jeśli nie naruszają one zasady proporcjonalności, są one uzasadnionym ograniczeniem wolności oskarżonego a nie reperkusją jego (jeszcze nie udowodnionej) winy. Nie chroni też ono oskarżonego przed koniecznością dowodzenia wysuwanych tez w myśl zasady \textit{ei incubit probatio qui dicit, non qui negat}, z oczywistym wyjątkiem samej niewinności (a więc brak jej dowodzenia i milczenie nie mogą mieć negatywnych konsekwencji).\medskip

Zasadą wynikającą z zasady domniemania niewinności jest wyrażona w art. 5 § 2 \acr{kpk} reguła \textit{in dubio pro reo}, której zastosowanie rozciąga sie zarówno na wątpliwości faktyczne jak i prawne i opiera się na nasępujących przesłankach:\begin{enumerate}
    \item wątpliwości nie da się usunąć;
    \item są to wątpliwości sądu (a nie stron).
\end{enumerate}

Reguła \textit{in dubio pro reo} napotyka \textit{de lege lata} ograniczenie w postaci instytucji wezwania przez sąd oskarżyciela publicznego do uzupełnienia dowodów, która, przed zastosowaniem tej reguły, nakłada na sędziów jeszcze jeden obowiązek.
